\documentclass{article}
\usepackage{amsmath, amssymb, geometry, enumitem}
\geometry{a4paper, margin=1in}

\title{Coordinate Systems: Comprehensive Problem Sets}
\author{}
\date{}

\begin{document}

\maketitle
\tableofcontents
\newpage

% ---------------------------------------------------------
\section{Concept Checklist \& Question Mapping}
% ---------------------------------------------------------
This section outlines all the core concepts, problem types, tricks, and pitfalls extracted from the provided text on Coordinate Systems (Section 4.4). Each concept is mapped to the corresponding questions in the problem set to ensure comprehensive coverage.

\begin{itemize}
    \item \textbf{Concept 1: Vector Reconstruction (Type 1)} 
    \newline \textit{Description:} Given a basis $B$ and the coordinate vector $[x]_B$, find the standard vector $x$ using $x = P_B[x]_B$.
    \newline \textit{Mapped Questions:} 1, 2, 3, 4
    
    \item \textbf{Concept 2: Coordinate Extraction via Row Reduction (Type 2)}
    \newline \textit{Description:} Given a vector $x$ and a basis $B$, find $[x]_B$ by setting up the augmented matrix $[P_B \mid x]$ and row reducing to $[I \mid [x]_B]$.
    \newline \textit{Mapped Questions:} 5, 6, 7, 8
    
    \item \textbf{Concept 3: Coordinate Extraction via Matrix Inverse (Type 3)}
    \newline \textit{Description:} Use the formula $[x]_B = P_B^{-1}x$. Highly efficient for $2 \times 2$ matrices using the shortcut formula.
    \newline \textit{Mapped Questions:} 9, 10, 11, 12
    
    \item \textbf{Concept 4: Change-of-Coordinates Matrix Construction}
    \newline \textit{Description:} Constructing $P_B$ by placing the basis vectors of $B$ as columns.
    \newline \textit{Mapped Questions:} 13, 14
    
    \item \textbf{Concept 5: Isomorphic Mapping for Abstract Spaces - Polynomial Coordinates (Type 4)}
    \newline \textit{Description:} Translating polynomials into column vectors in $\mathbb{R}^n$ by extracting coefficients relative to a standard basis.
    \newline \textit{Mapped Questions:} 15, 16, 17, 18
    
    \item \textbf{Concept 6: Isomorphic Mapping - Linear Independence (Type 4)}
    \newline \textit{Description:} Testing linear independence of abstract objects (like polynomials) by mapping them to $\mathbb{R}^n$ and checking for pivot columns.
    \newline \textit{Mapped Questions:} 19, 20, 21
    
    \item \textbf{Concept 7: Change of Basis Between Two Non-Standard Bases (Part 7 Variation)}
    \newline \textit{Description:} Finding $P_{C \leftarrow B}$ by evaluating the $C$-coordinates of each basis vector in $B$.
    \newline \textit{Mapped Questions:} 22, 23, 24
    
    \item \textbf{Concept 8: Conceptual Synthesis \& True/False}
    \newline \textit{Description:} Testing the theoretical underpinnings of isomorphisms, basis definitions, and standard bases.
    \newline \textit{Mapped Questions:} 25, 26, 27
    
    \item \textbf{Concept 9: Diagnostic Testing - "Find the Flaw" (Part 8 Pitfalls)}
    \newline \textit{Description:} Identifying common mistakes such as reversing the mapping ($P_B$ vs $P_B^{-1}$), placing vectors as rows instead of columns, and forgetting the determinant in the $2 \times 2$ inverse.
    \newline \textit{Mapped Questions:} 28, 29, 30
\end{itemize}

\newpage

% ---------------------------------------------------------
\section{Problem Sets}
% ---------------------------------------------------------

\subsection*{Type 1: Vector Reconstruction}
\textbf{1.} Let $B = \{b_1, b_2\}$ be a basis for $\mathbb{R}^2$ where $b_1 = \begin{bmatrix} 1 \\ 2 \end{bmatrix}$ and $b_2 = \begin{bmatrix} 3 \\ 4 \end{bmatrix}$. Find the vector $x$ determined by the coordinate vector $[x]_B = \begin{bmatrix} 2 \\ -1 \end{bmatrix}$.

\vspace{0.2cm}
\textbf{2.} Let $B = \{b_1, b_2\}$ be a basis for $\mathbb{R}^2$ where $b_1 = \begin{bmatrix} -1 \\ 3 \end{bmatrix}$ and $b_2 = \begin{bmatrix} 2 \\ -1 \end{bmatrix}$. Find the vector $x$ given $[x]_B = \begin{bmatrix} 4 \\ 2 \end{bmatrix}$.

\vspace{0.2cm}
\textbf{3.} Let $B = \{b_1, b_2, b_3\}$ be a basis for $\mathbb{R}^3$ where $b_1 = \begin{bmatrix} 1 \\ 0 \\ 1 \end{bmatrix}$, $b_2 = \begin{bmatrix} 0 \\ 1 \\ 1 \end{bmatrix}$, and $b_3 = \begin{bmatrix} 1 \\ 1 \\ 0 \end{bmatrix}$. Find $x$ if $[x]_B = \begin{bmatrix} 1 \\ 2 \\ -1 \end{bmatrix}$.

\vspace{0.2cm}
\textbf{4.} Given the basis $B = \left\{ \begin{bmatrix} 2 \\ -1 \\ 0 \end{bmatrix}, \begin{bmatrix} 0 \\ 2 \\ 1 \end{bmatrix}, \begin{bmatrix} 1 \\ 0 \\ -1 \end{bmatrix} \right\}$ and $[x]_B = \begin{bmatrix} -2 \\ 1 \\ 3 \end{bmatrix}$, find the standard vector $x$.

\subsection*{Type 2: Coordinate Extraction via Row Reduction}
\textbf{5.} Find the coordinate vector $[x]_B$ of $x$ relative to the basis $B = \{b_1, b_2, b_3\}$, where $b_1 = \begin{bmatrix} 1 \\ 1 \\ 1 \end{bmatrix}$, $b_2 = \begin{bmatrix} 0 \\ 1 \\ 1 \end{bmatrix}$, $b_3 = \begin{bmatrix} 0 \\ 0 \\ 1 \end{bmatrix}$, and $x = \begin{bmatrix} 2 \\ 5 \\ 8 \end{bmatrix}$.

\vspace{0.2cm}
\textbf{6.} Let $B = \left\{ \begin{bmatrix} 1 \\ 0 \\ -1 \end{bmatrix}, \begin{bmatrix} 2 \\ 1 \\ 0 \end{bmatrix}, \begin{bmatrix} 0 \\ 1 \\ 1 \end{bmatrix} \right\}$ and $x = \begin{bmatrix} 3 \\ 2 \\ 0 \end{bmatrix}$. Find $[x]_B$ using row reduction.

\vspace{0.2cm}
\textbf{7.} Find $[x]_B$ given $B = \left\{ \begin{bmatrix} 1 \\ 2 \\ 3 \end{bmatrix}, \begin{bmatrix} 0 \\ 1 \\ 2 \end{bmatrix}, \begin{bmatrix} 0 \\ 0 \\ 1 \end{bmatrix} \right\}$ and $x = \begin{bmatrix} 4 \\ 7 \\ 9 \end{bmatrix}$.

\vspace{0.2cm}
\textbf{8.} Let $B = \left\{ \begin{bmatrix} 2 \\ 0 \\ 0 \end{bmatrix}, \begin{bmatrix} 1 \\ 3 \\ 0 \end{bmatrix}, \begin{bmatrix} 1 \\ 1 \\ 1 \end{bmatrix} \right\}$ and $x = \begin{bmatrix} 5 \\ 4 \\ 1 \end{bmatrix}$. Find $[x]_B$.

\subsection*{Type 3: Coordinate Extraction via Matrix Inverse}
\textbf{9.} Use an inverse matrix to find $[x]_B$ for $B = \left\{ \begin{bmatrix} 3 \\ 1 \end{bmatrix}, \begin{bmatrix} 5 \\ 2 \end{bmatrix} \right\}$ and $x = \begin{bmatrix} 2 \\ -1 \end{bmatrix}$.

\vspace{0.2cm}
\textbf{10.} Use the $2 \times 2$ inverse shortcut to find $[x]_B$ for $B = \left\{ \begin{bmatrix} 4 \\ -2 \end{bmatrix}, \begin{bmatrix} -1 \\ 1 \end{bmatrix} \right\}$ and $x = \begin{bmatrix} 6 \\ 0 \end{bmatrix}$.

\vspace{0.2cm}
\textbf{11.} Find $[x]_B$ for $B = \left\{ \begin{bmatrix} 2 \\ 3 \end{bmatrix}, \begin{bmatrix} 1 \\ 4 \end{bmatrix} \right\}$ and $x = \begin{bmatrix} 7 \\ 13 \end{bmatrix}$.

\vspace{0.2cm}
\textbf{12.} Find $[x]_B$ for $B = \left\{ \begin{bmatrix} 5 \\ -3 \end{bmatrix}, \begin{bmatrix} -4 \\ 2 \end{bmatrix} \right\}$ and $x = \begin{bmatrix} -2 \\ 2 \end{bmatrix}$.

\subsection*{Type 4: Change-of-Coordinates Matrix Construction}
\textbf{13.} Find the change-of-coordinates matrix $P_B$ from $B$ to the standard basis in $\mathbb{R}^3$, given $B = \left\{ \begin{bmatrix} 1 \\ 2 \\ 3 \end{bmatrix}, \begin{bmatrix} 4 \\ 5 \\ 6 \end{bmatrix}, \begin{bmatrix} 7 \\ 8 \\ 9 \end{bmatrix} \right\}$.

\vspace{0.2cm}
\textbf{14.} Find the change-of-coordinates matrix $P_B$ for the basis $B = \left\{ \begin{bmatrix} -1 \\ 0 \end{bmatrix}, \begin{bmatrix} 0 \\ -1 \end{bmatrix} \right\}$.

\subsection*{Type 5: Isomorphic Mapping for Abstract Spaces (Polynomials)}
\textbf{15.} The set $B = \{1, t, t^2\}$ is the standard basis for $\mathbb{P}_2$. Find the coordinate vector of $p(t) = 3 - 2t + 4t^2$ relative to $B$.

\vspace{0.2cm}
\textbf{16.} The set $B = \{1+t, 1-t, t^2\}$ is a basis for $\mathbb{P}_2$. Find the coordinate vector of $p(t) = 4 + 2t - 3t^2$ relative to $B$.

\vspace{0.2cm}
\textbf{17.} Let $B = \{1, t+1, t^2+t+1\}$ be a basis for $\mathbb{P}_2$. Find $[p]_B$ for $p(t) = 2 - t + 3t^2$.

\vspace{0.2cm}
\textbf{18.} \textit{(Pitfall Warning)} Find the coordinate vector of $p(t) = 5t^2 - 3t + 2$ relative to the basis $B = \{t^2, t, 1\}$.

\subsection*{Type 6: Linear Independence via Isomorphism}
\textbf{19.} Use coordinate vectors to test whether the set of polynomials $\{1+t, t+t^2, 1+t^2\}$ is linearly independent in $\mathbb{P}_2$.

\vspace{0.2cm}
\textbf{20.} Test whether the polynomials $\{2-t, 1+t^2, 3-t+t^2\}$ are linearly independent.

\vspace{0.2cm}
\textbf{21.} Use coordinate vectors to test whether the set $\{1, 1+t, 1+t+t^2, 1+t+t^2+t^3\}$ is linearly independent in $\mathbb{P}_3$.

\subsection*{Type 7: Change of Basis Between Two Non-Standard Bases}
\textbf{22.} Let $B = \left\{ \begin{bmatrix} 1 \\ 1 \end{bmatrix}, \begin{bmatrix} 0 \\ 1 \end{bmatrix} \right\}$ and $C = \left\{ \begin{bmatrix} 1 \\ 0 \end{bmatrix}, \begin{bmatrix} 1 \\ 1 \end{bmatrix} \right\}$. Find the change-of-coordinates matrix from $B$ to $C$, denoted $P_{C \leftarrow B}$.

\vspace{0.2cm}
\textbf{23.} Let $B = \left\{ \begin{bmatrix} 2 \\ 1 \end{bmatrix}, \begin{bmatrix} 1 \\ 2 \end{bmatrix} \right\}$ and $C = \left\{ \begin{bmatrix} 1 \\ 1 \end{bmatrix}, \begin{bmatrix} 1 \\ -1 \end{bmatrix} \right\}$. Find $P_{C \leftarrow B}$.

\vspace{0.2cm}
\textbf{24.} Let $B = \left\{ \begin{bmatrix} 1 \\ 0 \\ 0 \end{bmatrix}, \begin{bmatrix} 1 \\ 1 \\ 0 \end{bmatrix}, \begin{bmatrix} 1 \\ 1 \\ 1 \end{bmatrix} \right\}$ and $C = \left\{ \begin{bmatrix} 0 \\ 0 \\ 1 \end{bmatrix}, \begin{bmatrix} 0 \\ 1 \\ 0 \end{bmatrix}, \begin{bmatrix} 1 \\ 0 \\ 0 \end{bmatrix} \right\}$. Find $P_{C \leftarrow B}$.

\subsection*{Type 8: Conceptual Synthesis (True/False)}
\textbf{25.} (T/F) The columns of the change-of-coordinates matrix $P_B$ are linearly dependent. Justify your answer.

\vspace{0.2cm}
\textbf{26.} (T/F) If $V$ is a vector space of dimension $n$, then $V$ is isomorphic to $\mathbb{R}^n$. Justify your answer.

\vspace{0.2cm}
\textbf{27.} (T/F) The coordinate mapping $x \mapsto [x]_B$ is a linear transformation. Justify your answer.

\subsection*{Type 9: Diagnostic Testing - "Find the Flaw"}
\textbf{28.} \textbf{Question:} Given $B = \left\{ \begin{bmatrix} 2 \\ 1 \end{bmatrix}, \begin{bmatrix} 1 \\ 1 \end{bmatrix} \right\}$ and $x = \begin{bmatrix} 3 \\ 2 \end{bmatrix}$, find $[x]_B$. 
\newline \textbf{Flawed Solution:} $[x]_B = P_B x = \begin{bmatrix} 2 & 1 \\ 1 & 1 \end{bmatrix} \begin{bmatrix} 3 \\ 2 \end{bmatrix} = \begin{bmatrix} 8 \\ 5 \end{bmatrix}$.
\newline Identify the flaw and provide the correct solution.

\vspace{0.2cm}
\textbf{29.} \textbf{Question:} Are the polynomials $p_1 = 1+t$ and $p_2 = 1-t$ linearly independent?
\newline \textbf{Flawed Solution:} We map them to $\mathbb{R}^2$ and place them as rows in a matrix: $A = \begin{bmatrix} 1 & 1 \\ 1 & -1 \end{bmatrix} \sim \begin{bmatrix} 1 & 1 \\ 0 & -2 \end{bmatrix}$. Since there are 2 pivots, they are independent.
\newline Identify the conceptual flaw in the setup, even if the final conclusion is accidentally correct.

\vspace{0.2cm}
\textbf{30.} \textbf{Question:} Find the inverse of the change-of-coordinates matrix $P_B = \begin{bmatrix} 4 & 2 \\ 5 & 3 \end{bmatrix}$.
\newline \textbf{Flawed Solution:} Using the $2 \times 2$ shortcut, we swap the main diagonal and negate the anti-diagonal: $P_B^{-1} = \begin{bmatrix} 3 & -2 \\ -5 & 4 \end{bmatrix}$.
\newline Identify the flaw and provide the correct solution.

\newpage

% ---------------------------------------------------------
\section{Detailed Solutions}
% ---------------------------------------------------------

\textbf{Solution 1:} By definition, $x = P_B[x]_B = c_1b_1 + c_2b_2$.
\[ x = 2 \begin{bmatrix} 1 \\ 2 \end{bmatrix} + (-1) \begin{bmatrix} 3 \\ 4 \end{bmatrix} = \begin{bmatrix} 2 \\ 4 \end{bmatrix} + \begin{bmatrix} -3 \\ -4 \end{bmatrix} = \begin{bmatrix} -1 \\ 0 \end{bmatrix} \]
\textbf{Final Answer:} $x = \begin{bmatrix} -1 \\ 0 \end{bmatrix}$

\vspace{0.3cm}
\textbf{Solution 2:} $x = c_1b_1 + c_2b_2$.
\[ x = 4 \begin{bmatrix} -1 \\ 3 \end{bmatrix} + 2 \begin{bmatrix} 2 \\ -1 \end{bmatrix} = \begin{bmatrix} -4 \\ 12 \end{bmatrix} + \begin{bmatrix} 4 \\ -2 \end{bmatrix} = \begin{bmatrix} 0 \\ 10 \end{bmatrix} \]
\textbf{Final Answer:} $x = \begin{bmatrix} 0 \\ 10 \end{bmatrix}$

\vspace{0.3cm}
\textbf{Solution 3:} $x = c_1b_1 + c_2b_2 + c_3b_3$.
\[ x = 1 \begin{bmatrix} 1 \\ 0 \\ 1 \end{bmatrix} + 2 \begin{bmatrix} 0 \\ 1 \\ 1 \end{bmatrix} - 1 \begin{bmatrix} 1 \\ 1 \\ 0 \end{bmatrix} = \begin{bmatrix} 1 \\ 0 \\ 1 \end{bmatrix} + \begin{bmatrix} 0 \\ 2 \\ 2 \end{bmatrix} + \begin{bmatrix} -1 \\ -1 \\ 0 \end{bmatrix} = \begin{bmatrix} 0 \\ 1 \\ 3 \end{bmatrix} \]
\textbf{Final Answer:} $x = \begin{bmatrix} 0 \\ 1 \\ 3 \end{bmatrix}$

\vspace{0.3cm}
\textbf{Solution 4:} $x = c_1b_1 + c_2b_2 + c_3b_3$.
\[ x = -2 \begin{bmatrix} 2 \\ -1 \\ 0 \end{bmatrix} + 1 \begin{bmatrix} 0 \\ 2 \\ 1 \end{bmatrix} + 3 \begin{bmatrix} 1 \\ 0 \\ -1 \end{bmatrix} = \begin{bmatrix} -4 \\ 2 \\ 0 \end{bmatrix} + \begin{bmatrix} 0 \\ 2 \\ 1 \end{bmatrix} + \begin{bmatrix} 3 \\ 0 \\ -3 \end{bmatrix} = \begin{bmatrix} -1 \\ 4 \\ -2 \end{bmatrix} \]
\textbf{Final Answer:} $x = \begin{bmatrix} -1 \\ 4 \\ -2 \end{bmatrix}$

\vspace{0.3cm}
\textbf{Solution 5:} Set up the augmented matrix $[P_B \mid x]$:
\[ \begin{bmatrix} 1 & 0 & 0 & \big| & 2 \\ 1 & 1 & 0 & \big| & 5 \\ 1 & 1 & 1 & \big| & 8 \end{bmatrix} \]
From $R_1$, $c_1 = 2$. Substitute into $R_2$: $2 + c_2 = 5 \implies c_2 = 3$. Substitute into $R_3$: $2 + 3 + c_3 = 8 \implies c_3 = 3$.
\textbf{Final Answer:} $[x]_B = \begin{bmatrix} 2 \\ 3 \\ 3 \end{bmatrix}$

\vspace{0.3cm}
\textbf{Solution 6:} Set up the augmented matrix $[P_B \mid x]$:
\[ \begin{bmatrix} 1 & 2 & 0 & \big| & 3 \\ 0 & 1 & 1 & \big| & 2 \\ -1 & 0 & 1 & \big| & 0 \end{bmatrix} \]
$R_3 \leftarrow R_3 + R_1$:
\[ \begin{bmatrix} 1 & 2 & 0 & \big| & 3 \\ 0 & 1 & 1 & \big| & 2 \\ 0 & 2 & 1 & \big| & 3 \end{bmatrix} \]
$R_3 \leftarrow R_3 - 2R_2$:
\[ \begin{bmatrix} 1 & 2 & 0 & \big| & 3 \\ 0 & 1 & 1 & \big| & 2 \\ 0 & 0 & -1 & \big| & -1 \end{bmatrix} \implies c_3 = 1 \]
Back-substitute: $c_2 + 1 = 2 \implies c_2 = 1$. $c_1 + 2(1) = 3 \implies c_1 = 1$.
\textbf{Final Answer:} $[x]_B = \begin{bmatrix} 1 \\ 1 \\ 1 \end{bmatrix}$

\vspace{0.3cm}
\textbf{Solution 7:} Set up the augmented matrix $[P_B \mid x]$:
\[ \begin{bmatrix} 1 & 0 & 0 & \big| & 4 \\ 2 & 1 & 0 & \big| & 7 \\ 3 & 2 & 1 & \big| & 9 \end{bmatrix} \]
$c_1 = 4$. $2(4) + c_2 = 7 \implies c_2 = -1$. $3(4) + 2(-1) + c_3 = 9 \implies 12 - 2 + c_3 = 9 \implies c_3 = -1$.
\textbf{Final Answer:} $[x]_B = \begin{bmatrix} 4 \\ -1 \\ -1 \end{bmatrix}$

\vspace{0.3cm}
\textbf{Solution 8:} Set up the augmented matrix $[P_B \mid x]$:
\[ \begin{bmatrix} 2 & 1 & 1 & \big| & 5 \\ 0 & 3 & 1 & \big| & 4 \\ 0 & 0 & 1 & \big| & 1 \end{bmatrix} \]
$c_3 = 1$. $3c_2 + 1 = 4 \implies 3c_2 = 3 \implies c_2 = 1$. $2c_1 + 1 + 1 = 5 \implies 2c_1 = 3 \implies c_1 = 1.5$.
\textbf{Final Answer:} $[x]_B = \begin{bmatrix} 1.5 \\ 1 \\ 1 \end{bmatrix}$

\vspace{0.3cm}
\textbf{Solution 9:} $P_B = \begin{bmatrix} 3 & 5 \\ 1 & 2 \end{bmatrix}$. $\det(P_B) = (3)(2) - (5)(1) = 1$.
$P_B^{-1} = \begin{bmatrix} 2 & -5 \\ -1 & 3 \end{bmatrix}$.
$[x]_B = P_B^{-1}x = \begin{bmatrix} 2 & -5 \\ -1 & 3 \end{bmatrix} \begin{bmatrix} 2 \\ -1 \end{bmatrix} = \begin{bmatrix} 4 + 5 \\ -2 - 3 \end{bmatrix} = \begin{bmatrix} 9 \\ -5 \end{bmatrix}$.
\textbf{Final Answer:} $[x]_B = \begin{bmatrix} 9 \\ -5 \end{bmatrix}$

\vspace{0.3cm}
\textbf{Solution 10:} $P_B = \begin{bmatrix} 4 & -1 \\ -2 & 1 \end{bmatrix}$. $\det(P_B) = (4)(1) - (-1)(-2) = 4 - 2 = 2$.
$P_B^{-1} = \frac{1}{2} \begin{bmatrix} 1 & 1 \\ 2 & 4 \end{bmatrix}$.
$[x]_B = \frac{1}{2} \begin{bmatrix} 1 & 1 \\ 2 & 4 \end{bmatrix} \begin{bmatrix} 6 \\ 0 \end{bmatrix} = \frac{1}{2} \begin{bmatrix} 6 \\ 12 \end{bmatrix} = \begin{bmatrix} 3 \\ 6 \end{bmatrix}$.
\textbf{Final Answer:} $[x]_B = \begin{bmatrix} 3 \\ 6 \end{bmatrix}$

\vspace{0.3cm}
\textbf{Solution 11:} $P_B = \begin{bmatrix} 2 & 1 \\ 3 & 4 \end{bmatrix}$. $\det(P_B) = 8 - 3 = 5$.
$P_B^{-1} = \frac{1}{5} \begin{bmatrix} 4 & -1 \\ -3 & 2 \end{bmatrix}$.
$[x]_B = \frac{1}{5} \begin{bmatrix} 4 & -1 \\ -3 & 2 \end{bmatrix} \begin{bmatrix} 7 \\ 13 \end{bmatrix} = \frac{1}{5} \begin{bmatrix} 28 - 13 \\ -21 + 26 \end{bmatrix} = \frac{1}{5} \begin{bmatrix} 15 \\ 5 \end{bmatrix} = \begin{bmatrix} 3 \\ 1 \end{bmatrix}$.
\textbf{Final Answer:} $[x]_B = \begin{bmatrix} 3 \\ 1 \end{bmatrix}$

\vspace{0.3cm}
\textbf{Solution 12:} $P_B = \begin{bmatrix} 5 & -4 \\ -3 & 2 \end{bmatrix}$. $\det(P_B) = 10 - 12 = -2$.
$P_B^{-1} = -\frac{1}{2} \begin{bmatrix} 2 & 4 \\ 3 & 5 \end{bmatrix}$.
$[x]_B = -\frac{1}{2} \begin{bmatrix} 2 & 4 \\ 3 & 5 \end{bmatrix} \begin{bmatrix} -2 \\ 2 \end{bmatrix} = -\frac{1}{2} \begin{bmatrix} -4 + 8 \\ -6 + 10 \end{bmatrix} = -\frac{1}{2} \begin{bmatrix} 4 \\ 4 \end{bmatrix} = \begin{bmatrix} -2 \\ -2 \end{bmatrix}$.
\textbf{Final Answer:} $[x]_B = \begin{bmatrix} -2 \\ -2 \end{bmatrix}$

\vspace{0.3cm}
\textbf{Solution 13:} The change-of-coordinates matrix is formed by placing the basis vectors as columns.
\textbf{Final Answer:} $P_B = \begin{bmatrix} 1 & 4 & 7 \\ 2 & 5 & 8 \\ 3 & 6 & 9 \end{bmatrix}$

\vspace{0.3cm}
\textbf{Solution 14:} Place the basis vectors as columns.
\textbf{Final Answer:} $P_B = \begin{bmatrix} -1 & 0 \\ 0 & -1 \end{bmatrix}$

\vspace{0.3cm}
\textbf{Solution 15:} The basis is ordered $1, t, t^2$. The coefficients of $p(t) = 3 - 2t + 4t^2$ correspond directly to this order.
\textbf{Final Answer:} $[p]_B = \begin{bmatrix} 3 \\ -2 \\ 4 \end{bmatrix}$

\vspace{0.3cm}
\textbf{Solution 16:} We need $c_1(1+t) + c_2(1-t) + c_3(t^2) = 4 + 2t - 3t^2$.
Group by powers of $t$: $(c_1 + c_2) + (c_1 - c_2)t + c_3t^2 = 4 + 2t - 3t^2$.
Equating coefficients:
$c_1 + c_2 = 4$
$c_1 - c_2 = 2$
$c_3 = -3$
Adding the first two equations: $2c_1 = 6 \implies c_1 = 3$. Then $3 + c_2 = 4 \implies c_2 = 1$.
\textbf{Final Answer:} $[p]_B = \begin{bmatrix} 3 \\ 1 \\ -3 \end{bmatrix}$

\vspace{0.3cm}
\textbf{Solution 17:} We need $c_1(1) + c_2(t+1) + c_3(t^2+t+1) = 2 - t + 3t^2$.
Group by powers of $t$: $(c_1 + c_2 + c_3) + (c_2 + c_3)t + c_3t^2 = 2 - t + 3t^2$.
Equating coefficients:
$c_3 = 3$
$c_2 + c_3 = -1 \implies c_2 + 3 = -1 \implies c_2 = -4$
$c_1 + c_2 + c_3 = 2 \implies c_1 - 4 + 3 = 2 \implies c_1 - 1 = 2 \implies c_1 = 3$.
\textbf{Final Answer:} $[p]_B = \begin{bmatrix} 3 \\ -4 \\ 3 \end{bmatrix}$

\vspace{0.3cm}
\textbf{Solution 18:} The basis is ordered $t^2, t, 1$. We must extract the coefficients of $p(t) = 5t^2 - 3t + 2$ in this exact order.
\textbf{Final Answer:} $[p]_B = \begin{bmatrix} 5 \\ -3 \\ 2 \end{bmatrix}$

\vspace{0.3cm}
\textbf{Solution 19:} Map the polynomials to $\mathbb{R}^3$ using the standard basis $\{1, t, t^2\}$:
$[p_1]_E = \begin{bmatrix} 1 \\ 1 \\ 0 \end{bmatrix}, [p_2]_E = \begin{bmatrix} 0 \\ 1 \\ 1 \end{bmatrix}, [p_3]_E = \begin{bmatrix} 1 \\ 0 \\ 1 \end{bmatrix}$.
Place as columns in a matrix: $A = \begin{bmatrix} 1 & 0 & 1 \\ 1 & 1 & 0 \\ 0 & 1 & 1 \end{bmatrix}$.
Row reduce or find determinant: $\det(A) = 1(1-0) - 0 + 1(1-0) = 2 \neq 0$.
Since the determinant is non-zero, the columns are linearly independent.
\textbf{Final Answer:} The polynomials are linearly independent.

\vspace{0.3cm}
\textbf{Solution 20:} Map to $\mathbb{R}^3$: $[p_1]_E = \begin{bmatrix} 2 \\ -1 \\ 0 \end{bmatrix}, [p_2]_E = \begin{bmatrix} 1 \\ 0 \\ 1 \end{bmatrix}, [p_3]_E = \begin{bmatrix} 3 \\ -1 \\ 1 \end{bmatrix}$.
Notice that $p_1 + p_2 = (2-t) + (1+t^2) = 3 - t + t^2 = p_3$.
Because there is a non-trivial linear dependence relation ($p_1 + p_2 - p_3 = 0$), they are dependent.
\textbf{Final Answer:} The polynomials are linearly dependent.

\vspace{0.3cm}
\textbf{Solution 21:} Map to $\mathbb{R}^4$ using $\{1, t, t^2, t^3\}$:
$A = \begin{bmatrix} 1 & 1 & 1 & 1 \\ 0 & 1 & 1 & 1 \\ 0 & 0 & 1 & 1 \\ 0 & 0 & 0 & 1 \end{bmatrix}$.
The matrix is already in upper echelon form with a pivot in every column.
\textbf{Final Answer:} The polynomials are linearly independent.

\vspace{0.3cm}
\textbf{Solution 22:} We need to find $[b_1]_C$ and $[b_2]_C$.
For $[b_1]_C$: $c_1 \begin{bmatrix} 1 \\ 0 \end{bmatrix} + c_2 \begin{bmatrix} 1 \\ 1 \end{bmatrix} = \begin{bmatrix} 1 \\ 1 \end{bmatrix} \implies c_2 = 1, c_1 = 0$. So $[b_1]_C = \begin{bmatrix} 0 \\ 1 \end{bmatrix}$.
For $[b_2]_C$: $c_1 \begin{bmatrix} 1 \\ 0 \end{bmatrix} + c_2 \begin{bmatrix} 1 \\ 1 \end{bmatrix} = \begin{bmatrix} 0 \\ 1 \end{bmatrix} \implies c_2 = 1, c_1 + 1 = 0 \implies c_1 = -1$. So $[b_2]_C = \begin{bmatrix} -1 \\ 1 \end{bmatrix}$.
\textbf{Final Answer:} $P_{C \leftarrow B} = \begin{bmatrix} 0 & -1 \\ 1 & 1 \end{bmatrix}$

\vspace{0.3cm}
\textbf{Solution 23:}
For $[b_1]_C$: $c_1 \begin{bmatrix} 1 \\ 1 \end{bmatrix} + c_2 \begin{bmatrix} 1 \\ -1 \end{bmatrix} = \begin{bmatrix} 2 \\ 1 \end{bmatrix} \implies c_1+c_2=2, c_1-c_2=1$. Adding gives $2c_1=3 \implies c_1=1.5$. Then $c_2=0.5$. $[b_1]_C = \begin{bmatrix} 1.5 \\ 0.5 \end{bmatrix}$.
For $[b_2]_C$: $c_1 \begin{bmatrix} 1 \\ 1 \end{bmatrix} + c_2 \begin{bmatrix} 1 \\ -1 \end{bmatrix} = \begin{bmatrix} 1 \\ 2 \end{bmatrix} \implies c_1+c_2=1, c_1-c_2=2$. Adding gives $2c_1=3 \implies c_1=1.5$. Then $c_2=-0.5$. $[b_2]_C = \begin{bmatrix} 1.5 \\ -0.5 \end{bmatrix}$.
\textbf{Final Answer:} $P_{C \leftarrow B} = \begin{bmatrix} 1.5 & 1.5 \\ 0.5 & -0.5 \end{bmatrix}$

\vspace{0.3cm}
\textbf{Solution 24:} Notice that $C$ is just a reordering of the standard basis $E = \{e_3, e_2, e_1\}$.
$b_1 = \begin{bmatrix} 1 \\ 0 \\ 0 \end{bmatrix} = 0 \cdot e_3 + 0 \cdot e_2 + 1 \cdot e_1 \implies [b_1]_C = \begin{bmatrix} 0 \\ 0 \\ 1 \end{bmatrix}$.
$b_2 = \begin{bmatrix} 1 \\ 1 \\ 0 \end{bmatrix} = 0 \cdot e_3 + 1 \cdot e_2 + 1 \cdot e_1 \implies [b_2]_C = \begin{bmatrix} 0 \\ 1 \\ 1 \end{bmatrix}$.
$b_3 = \begin{bmatrix} 1 \\ 1 \\ 1 \end{bmatrix} = 1 \cdot e_3 + 1 \cdot e_2 + 1 \cdot e_1 \implies [b_3]_C = \begin{bmatrix} 1 \\ 1 \\ 1 \end{bmatrix}$.
\textbf{Final Answer:} $P_{C \leftarrow B} = \begin{bmatrix} 0 & 0 & 1 \\ 0 & 1 & 1 \\ 1 & 1 & 1 \end{bmatrix}$

\vspace{0.3cm}
\textbf{Solution 25:} \textbf{False.} By definition, a basis $B$ consists of linearly independent vectors. Since $P_B$ is formed by placing these basis vectors as its columns, the columns of $P_B$ are strictly linearly independent.

\vspace{0.3cm}
\textbf{Solution 26:} \textbf{True.} The coordinate mapping $x \mapsto [x]_B$ is a one-to-one and onto linear transformation (an isomorphism) from $V$ to $\mathbb{R}^n$. This means they share the exact same structural properties.

\vspace{0.3cm}
\textbf{Solution 27:} \textbf{True.} The coordinate mapping preserves vector addition ($[u+v]_B = [u]_B + [v]_B$) and scalar multiplication ($[cu]_B = c[u]_B$), which are the two requirements for a linear transformation.

\vspace{0.3cm}
\textbf{Solution 28:} \textbf{Flaw:} The student multiplied $x$ by $P_B$ instead of $P_B^{-1}$. $P_B$ maps coordinates to vectors ($x = P_B[x]_B$), not vectors to coordinates.
\textbf{Correct Step:} $[x]_B = P_B^{-1}x$.
$P_B^{-1} = \frac{1}{2-1} \begin{bmatrix} 1 & -1 \\ -1 & 2 \end{bmatrix} = \begin{bmatrix} 1 & -1 \\ -1 & 2 \end{bmatrix}$.
$[x]_B = \begin{bmatrix} 1 & -1 \\ -1 & 2 \end{bmatrix} \begin{bmatrix} 3 \\ 2 \end{bmatrix} = \begin{bmatrix} 3 - 2 \\ -3 + 4 \end{bmatrix} = \begin{bmatrix} 1 \\ 1 \end{bmatrix}$.

\vspace{0.3cm}
\textbf{Solution 29:} \textbf{Flaw:} The student placed the coordinate vectors as \textit{rows} instead of \textit{columns}. While row rank equals column rank (meaning the pivot count will accidentally yield the correct True/False conclusion for independence), placing them as rows destroys the ability to find the linear dependence relation $c_1p_1 + c_2p_2 = 0$. Coordinate vectors must always be placed as columns to solve $Ax = 0$.

\vspace{0.3cm}
\textbf{Solution 30:} \textbf{Flaw:} The student forgot to multiply the matrix by the scalar $\frac{1}{\det(P_B)}$.
\textbf{Correct Step:} $\det(P_B) = (4)(3) - (2)(5) = 12 - 10 = 2$.
$P_B^{-1} = \frac{1}{2} \begin{bmatrix} 3 & -2 \\ -5 & 4 \end{bmatrix} = \begin{bmatrix} 1.5 & -1 \\ -2.5 & 2 \end{bmatrix}$.

\newpage

% ---------------------------------------------------------
\section*{Concept Checklist \& Question Mapping (Review)}
% ---------------------------------------------------------
\begin{itemize}
    \item \textbf{Concept 1: Vector Reconstruction (Type 1)} $\rightarrow$ Q1, Q2, Q3, Q4
    \item \textbf{Concept 2: Coordinate Extraction via Row Reduction (Type 2)} $\rightarrow$ Q5, Q6, Q7, Q8
    \item \textbf{Concept 3: Coordinate Extraction via Matrix Inverse (Type 3)} $\rightarrow$ Q9, Q10, Q11, Q12
    \item \textbf{Concept 4: Change-of-Coordinates Matrix Construction} $\rightarrow$ Q13, Q14
    \item \textbf{Concept 5: Isomorphic Mapping for Abstract Spaces - Polynomial Coordinates (Type 4)} $\rightarrow$ Q15, Q16, Q17, Q18
    \item \textbf{Concept 6: Isomorphic Mapping - Linear Independence (Type 4)} $\rightarrow$ Q19, Q20, Q21
    \item \textbf{Concept 7: Change of Basis Between Two Non-Standard Bases (Part 7 Variation)} $\rightarrow$ Q22, Q23, Q24
    \item \textbf{Concept 8: Conceptual Synthesis \& True/False} $\rightarrow$ Q25, Q26, Q27
    \item \textbf{Concept 9: Diagnostic Testing - "Find the Flaw" (Part 8 Pitfalls)} $\rightarrow$ Q28, Q29, Q30
\end{itemize}

\end{document}